%\documentclass[12pt,oneside]{amsart}
\documentclass{scrartcl}

%\documentclass{article}

\usepackage{amsthm,amsmath,amssymb}
\usepackage[numbers]{natbib}

%\usepackage[colorlinks]{hyperref}
\usepackage{hyperref}
\hypersetup{
    colorlinks,%
    citecolor=black,%
    filecolor=black,%
    linkcolor=black,%
    urlcolor=black
}
\usepackage{hypernat}
\usepackage[top=2.5cm, bottom=2.5cm, left=2.8cm,
right=2.8cm]{geometry}
\usepackage{graphicx}


%\setlength{\parskip}{0pt}
%\setlength{\parsep}{0pt}
%\setlength{\headsep}{0pt}
%\setlength{\topskip}{0pt}
%\setlength{\topmargin}{0pt}
%\setlength{\topsep}{0pt}
%\setlength{\partopsep}{0pt}

\linespread{1}

%\usepackage{marvosym}

%\textwidth = 6.5 in \textheight = 9.2 in \oddsidemargin = 0.0 in
%\evensidemargin = 0.0 in \topmargin = -0.0 in \headheight = 0.0 in
%\headsep = 0.0in \parskip = 0.0in \parindent = 0.0in
%\def\baselinestretch{0.871}


\newtheorem{theorem}{Theorem}[section]
\newtheorem{proposition}[theorem]{Proposition}
\newtheorem{problem}[theorem]{Problem}
\newtheorem{fact}[theorem]{Fact}
\newtheorem{lemma}[theorem]{Lemma}
\newtheorem{defn}[theorem]{Definition}
\newtheorem{corollary}[theorem]{Corollary}
\newtheorem{remark}{Remark}[section]
\newtheorem{example}[theorem]{Example}

\newcommand{\E}{{\mathbb E}}
\newcommand{\se}{{\mathcal{E}}}
\newcommand{\R}{{\mathbb R}}
\renewcommand{\P}{{\mathbb P}}
\newcommand{\ac}{{\mathcal{A}}}
\newcommand{\A}{{\mathcal{A}}}
\newcommand{\B}{{\mathcal{B}}}
\newcommand{\C}{{\mathcal{C}}}
\newcommand{\M}{{\mathcal{M}}}
\newcommand{\Mf}{{\mathfrak{M}}}
\newcommand{\T}{{\mathcal{T}}}
\newcommand{\Z}{{\mathcal{Z}}}
\newcommand{\K}{{\mathcal{K}}}
\newcommand{\lc}{{\mathcal{L}}}
\newcommand{\Ps}{{\mathcal{P}}}
\newcommand{\G}{{\mathcal{G}}}
\newcommand{\pa}{{\mathring{p}}}
\newcommand{\F}{{\cal F}}
\newcommand{\samp}{{\mathcal{X}}}
\newcommand{\X}{{\mathcal{X}}}
\newcommand{\hi}{{\hat{\Phi}}}
\newcommand{\data}{{\text{data}}}
\newcommand{\relu}{{\text{ReLU}}}

\newcommand{\ridge}{{\hat{\beta}_{\text{ridge}}(\lambda)}}
\newcommand{\lasso}{{\hat{\beta}_{\text{lasso}}(\lambda)}}
\newcommand{\ridgemu}{{\hat{\mu}_t^{\text{ridge}}(\lambda)}}
\newcommand{\lassomu}{{\hat{\mu}_t^{\text{lasso}}(\lambda)}}


\newcounter{rcnt}[section]
\renewcommand{\thercnt}{(\roman{rcnt})}

\def\qt#1{\qquad\text{#1}}

\def\argmin{\mathop{\rm argmin}}
\def\argmax{\mathop{\rm argmax}}


\setlength{\parskip}{1.4 \medskipamount}
%\sloppy
%\linespread{1.3}

\begin{document}

\title{STAT 238 - Bayesian Statistics  \\
  Lecture Twenty Three} 
\subtitle{Spring 2026, UC Berkeley}

\author{Aditya Guntuboyina}


\date{18 March 2026}

\maketitle

\section{Last lecture: interpolation with Gaussian processes}
The interpolation problem is: Suppose we are given values of $f$ at
points $x_1, \dots, x_n$ in the domain $\Omega$ of $f$. Let $x \in
\Omega$ be a new point (i.e., distinct from  
$x_1, \dots, x_n$). What can we say about $f(x)$?

In the last lecture, we saw how to use Gaussian processes to solve
this problem. We model $\{f(x), x \in 
\Omega\}$ as a Gaussian process with mean zero and covariance
function or \textit{kernel} given by $K(x, x')$ i.e.,
\begin{equation*}
  \text{Cov}(f(x), f(x')) = K(x, x') ~~\text{for all $x, x' \in
    \Omega$}. 
\end{equation*}
Under this modeling assumption, the answer to the interpolation
question is given by the conditional distribution of $f(x)$ given
$f(x_1), \dots, f(x_n)$ which is calculated as follows.
Note that:
\begin{align*}
  (f(x_1), \dots, f(x_n), f(x)) \sim N \left(0, \begin{pmatrix} (K(x_i, 
                                             x_j))_{n \times n}  & (K(x_i,
                                                            x))_{n
                                                            \times 1}
                                             \\ (K(x, x_i))_{1
                                             \times n} & K(x,
                                                         x) \end{pmatrix}
  \right). 
\end{align*}
Using the notation $K = (K(x_i, 
                                             x_j))_{n \times n}$ and
                                             $\mathbf{k} = (K(x_i,
                                                            x))_{n
                                                            \times
                                                            1}$, we
                                                          can write
                                                          the
                                                          conditional
                                                          distribution
                                                          of
                                                          $f(x)$
                                                          given $f(x_1),
                                                          \dots, 
f(x_n)$ as
\begin{align*}
  f(x) \mid f(x_1), \dots, f(x_n) \sim N\left(\mathbf{k}^TK^{-1} (f(x_1), \dots,
  f(x_n))^T, K(x, x) -
  \mathbf{k}^T K^{-1} \mathbf{k} \right).
\end{align*}
Thus the posterior mean (or mode) estimate of $f(x)$ is given by
\begin{align}\label{posmean_int}
  \widehat{f(x)} = \mathbf{k}^TK^{-1} (f(x_1), \dots, f(x_n))^T. 
\end{align}

\section{Regression with Gaussian Processes}

Today, we will see how to perform regression using Gaussian processes.  
The key difference between \textbf{regression} and \textbf{interpolation} is that, in regression, the values \( f(x_1), \dots, f(x_n) \) are not observed exactly; instead, they are observed with noise.

More precisely, we are given observations \( y_1, \dots, y_n \) modeled as
\[
y_i = f(x_i) + \epsilon_i, \qquad \text{where } \epsilon_i \overset{\text{i.i.d.}}{\sim} N(0, \sigma^2).
\]
The parameter \( \sigma \) controls the level of noise, i.e., how much each observation \( y_i \) deviates from the true value \( f(x_i) \).

As in the interpolation setting, our goal is to estimate \( f(x) \) at a test point \( x \). This test point may be different from the observed inputs \( x_1, \dots, x_n \), or it may coincide with one of them.

Importantly, because the observations are noisy, it is meaningful to estimate \( f(x_i) \) even at observed points. In fact, by combining information from all observations, we can often obtain a better estimate of \( f(x_i) \) than the raw observation \( y_i \).

Another key difference from interpolation is that the input points \(
x_1, \dots, x_n \) need not necessarily be distinct. Since each observation
contains noise, having repeated measurements at the same input can
help improve the overall estimate.

The solution to the regression problem is very similar to that of the
interpolation problem with only one difference. The goal is to
estimate $f(x)$ at the test point $x$ given the data $(x_1, y_1),
\dots, (x_n, y_n)$. We will use the conditional distribution of $f(x)$
given $y_1, \dots, y_n$ (we will assume that $x_1, \dots, x_n, x$ are
deterministic). To calculate this conditional distribution, first note
that the marginal distribution of $(y_1, \dots, y_n, 
f(x_*))$ is given by
\begin{align*}
  (y_1, \dots, y_n, f(x)) \sim N \left(0, \begin{pmatrix} (K(x_i, 
                                             x_j))_{n \times n} +
                                             \sigma^2 I_n & (K(x_i,
                                                            x))_{n
                                                            \times 1}
                                             \\ (K(x, x_i))_{1
                                             \times n} & K(x,
                                                         x) \end{pmatrix}
  \right). 
\end{align*}
Using the notation $K = (K(x_i, 
                                             x_j))_{n \times n}$ and
                                             $\mathbf{k} = (K(x_i,
                                                            x))_{n
                                                            \times
                                                            1}$, we
                                                          can write
                                                          the
                                                          conditional
                                                          distribution
                                                          of
                                                          $f(x)$
                                                          given $y_1,
                                                          \dots, 
y_n$ as
\begin{align*}
  f(x) \mid \text{data} \sim N\left(\mathbf{k}^T \left(K +
  \sigma^2 I_n \right)^{-1} Y, K(x, x) - \mathbf{k}^T \left(K +
  \sigma^2 I_n
  \right)^{-1} \mathbf{k} \right).
\end{align*}
Thus the posterior mean (or mode) estimate of $f(x)$ is given by
\begin{align}\label{posmean}
  \widehat{f(x)} = \mathbf{k}^T \left(K +
  \sigma^2 I_n \right)^{-1} y, 
\end{align}
where $y$ is the $n \times 1$ vector with entries $y_1, \dots, y_n$.

So the only difference between \eqref{posmean} and \eqref{posmean_int}
is the presence of the additional $\sigma^2 I_n$ term in case of
regression.

Often, we would need to estimate $\sigma$ from the observed data (and
also additional hyperparameters present in the kernel $K$). For these,
the marginal likelihood of $y_1, \dots, y_n$ is important. This is
simply the multivariate normal distribution with mean vector 0 and
covariance $K + \sigma^2 I_n$. 

To illustrate the calculations, let us take the special case of the Integrated
Brownian Motion prior. 

\section{Calculations for the IBM prior}

We take the prior
\begin{align*}
  f(x) = \beta_0 + \beta_1 x + \tau I(x) 
\end{align*}
where $\beta_0, \beta_1$ are i.i.d $N(0, C)$ and $I(x)$ is integrated
Brownian Motion.

This is a Gaussian process prior with mean zero and covariance kernel:
\begin{align*}
  K(u, v) = C(1 + uv) + \tau^2 K_I(u, v)
\end{align*}
where $K_I(u, v)$ is the kernel corresponding to IBM, which is:
\begin{align*}
  K_I(u, v) &= \frac{1}{2} \left(\min(u, v) \right)^2\max(u, v)  -
              \frac{1}{6} \left(\min(u, v)\right)^3 \\
&= \frac{1}{6} \left(\min(u, v) \right)^2 \left(3 \max(u, v) - \min(u,
  v) \right) \\
&= u v (\min(u, v)) - \frac{u+v}{2} \left(\min(u, v) \right)^2 +
  \frac{1}{3} \left(\min(u, v) \right)^3. 
\end{align*}
Note that the kernel has the unknown parameter $\tau$ (which we write
as $\tau = \gamma \sigma$). The constant $C$ is assumed to be large
and it is not to be estimated (ideally we want to take $C =
+\infty$). 

Given data $(x_1, y_1), \dots, (x_n, y_n)$, what is the posterior of
$f(x)$? First let us assume that $\tau, \sigma$ are given. The
estimate is simply \eqref{posmean}. We simplify this expression
below. Let $X$ denote the $n \times 2$ matrix with columns $1$ and
$x_i$ (this is the usual $X$ matrix in the simple linear regression of
$y$ on $x$ based on the data $(x_1, y_1), \dots, (x_n, y_n)$). 

Note that
\begin{align*}
  \mathbf{k}^T &= (K(x, x_1), \dots, K(x, x_n)) \\
  &= \left(C(1 + x x_i) + \tau^2 K_I(x, x_i), i = 1, \dots, n \right)
  \\
               &= C (1, x) X^T + \tau^2 (K_I(x, x_1), \dots, K_I(x, x_n)) \\
  &= C (1, x) X^T + \tau^2 \mathbf{k}_I^T. 
\end{align*}
where
\begin{align*}
  \mathbf{k}_I^T := (K_I(x, x_1), \dots, K_I(x, x_n)),
\end{align*}
and $(1, x)$ denotes the row vector with entries $1$ and $x$. 

Further the $n \times n$ matrix $K$ has the $(i, j)$-th entry:
\begin{align*}
  C(1 + x_i x_j) + \gamma^2 \sigma^2 K_I(x_i, x_j)
\end{align*}
so that
\begin{align}\label{K_exp}
  K = C X X^T + \tau^2 K_I
\end{align}
where $K_I$ is the $n \times n$ matrix with $(i, j)$-th entry
$K_I(x_i, x_j)$. 

As a result, 
\begin{align*}
  \widehat{f(x)} =
  \mathbf{k}^T \left(K + \sigma^2 I_n \right)^{-1} y = \left(C (1, x)
  X^T + \tau^2 \mathbf{k}_I^T \right) \left(C X X^T + \tau^2 K_I +
  \sigma^2 I_n \right)^{-1} y. 
\end{align*}
This expression depends on the large constant $C$. Direct computation
with a large $C$ might make it unstable. It is therefore natural to
compute the limit as $C \rightarrow \infty$. Using the
Sherman-Morrison-Woodbury identity, 
\begin{align}
(K + \sigma^2 I_n)^{-1} &= (C X X^T + \tau^2 K_I + \sigma^2 I_n)^{-1}
                          \nonumber \\ &= \left(C X X^T + \Sigma
                                         \right)^{-1} \nonumber \\ &=
                                                           \Sigma^{-1}
                                                           -
                                                           \Sigma^{-1}
                                                           X
                                                           \left(C^{-1} 
                            I_2 + X^T \Sigma^{-1} X  \right)^{-1} X^T
                                                                                                                     \Sigma^{-1} \label{smw}
\end{align}
where
\begin{align}\label{Sigdef}
  \Sigma = \tau^2 K_I + \sigma^2 I_n. 
\end{align}
Further
\begin{align*}
  \mathbf{k} = C (1, x) X^T + \tau^2 \mathbf{k}_I^T. 
\end{align*}
We thus get
\begin{align*}
   \widehat{f(x)} &= \left(C (1, x) X^T + \tau^2 \mathbf{k}_I^T \right) \left(\Sigma^{-1} - \Sigma^{-1} X \left(C^{-1}
                            I_2 + X^T \Sigma^{-1} X  \right)^{-1} X^T
                \Sigma^{-1} \right) y. 
\end{align*}
Write $\tau = \gamma \sigma$. In Fact \ref{Climit}, it is proved that,
as $C \rightarrow \infty$ the above converges to:
\begin{align*}
   \widehat{f(x)} := (1, x) \left(X^T A_{\gamma}^{-1} X \right)^{-1} X^T
  A_{\gamma}^{-1} y + \gamma^2 \mathbf{k}_I^T \left(A_{\gamma}^{-1} -
  A_{\gamma}^{-1} X (X^T A_{\gamma}^{-1} X)^{-1} X^T A_{\gamma}^{-1}
  \right) y 
\end{align*}
where
\begin{align*}
    A_{\gamma} = I_{n \times n} + \gamma^2 K_I. 
\end{align*}
This expression for $\widehat{f(x)}$, which only depends on $\gamma = \tau/\sigma$
and not on $\tau$ and $\sigma$ individually, can be used for computation. 

\subsection{Estimation of $\gamma$ and $\sigma$}

The hyperparameters $\gamma$ and $\sigma$ also need to be estimated
from the observed data $(x_1, y_1), \dots, (x_n, y_n)$ (recall that
$\tau = \gamma \sigma$). For this, the marginal likelihood of the data
given $\sigma, \gamma$ is important. This is calculated using: 
\begin{align*}
  y \mid \sigma, \gamma \sim N(0, K + \sigma^2 I_n)
\end{align*}
where $K$  is given by \eqref{K_exp}. In other words,
\begin{align*}
  f_{y \mid \sigma, \gamma}(y) \propto \frac{1}{\sqrt{\det(K +
  \sigma^2 I_n)}} \exp \left(-\frac{1}{2} y^T (K + \sigma^2 I_n)^{-1}
  y \right). 
\end{align*}
For $(K + \sigma^2 I_n)^{-1}$, we use \eqref{smw} and let $C
\rightarrow \infty$ to get 
\begin{align*}
  (K + \sigma^2 I_n)^{-1} \rightarrow \Sigma^{-1} - \Sigma^{-1} X
  \left(X^T \Sigma^{-1} X  \right)^{-1} X^T \Sigma^{-1}.  
\end{align*}
Further, using the matrix determinant lemma, we get
\begin{align*}
  |K + \sigma^2 I_n| &= |\Sigma + C X X^T| \\
                     &= |\Sigma| |I + C X^T \Sigma^{-1} X| \\
                     &\approx |\Sigma| |C X^T \Sigma^{-1} X| \text{ when $C$ is large} \\
  &= |\Sigma| C^2 |X^T \Sigma^{-1} X| \propto |\Sigma| |X^T
    \Sigma^{-1} X|. 
\end{align*}
So the marginal likelihood of $y$ given $\tau, \sigma$ becomes:
\begin{align*}
  f_{y \mid \tau, \sigma}(y) \propto |\Sigma|^{-1/2} |X^T \Sigma^{-1}
  X|^{-1/2} \exp \left(-\frac{1}{2} y^T \left[\Sigma^{-1} -
  \Sigma^{-1} X \left(X^T \Sigma^{-1} X  \right)^{-1} X^T 
  \Sigma^{-1}\right] y  \right)
\end{align*}
with $\Sigma$ defined in \eqref{Sigdef}. 

We now take $\tau = \gamma \sigma$ so that
\begin{align*}
  \Sigma = \sigma^2 \left(I_n + \gamma^2 K_I \right) = \sigma^2
  A_{\gamma} ~~ \text{ where $A_{\gamma} := I_n + \gamma^2 K_I$}. 
\end{align*}
This gives
\begin{align*}
  f_{y \mid \gamma, \sigma}(y) \propto \sigma^{-(n-2)}
  |A_{\gamma}|^{-1/2} |X^T A_{\gamma}^{-1} X|^{-1/2} \exp
  \left(-\frac{y^T \left[A_{\gamma}^{-1} - A_{\gamma}^{-1} X (X^T
  A_{\gamma}^{-1} X)^{-1} X^T A_{\gamma}^{-1} \right] y}{2 \sigma^2}
  \right). 
\end{align*}
Combining this with the prior $\log \sigma \mid \gamma \sim
\text{uniform}(-\infty, \infty)$ gives
\begin{align*}
  \frac{1}{\sigma^2} \mid y, \gamma \sim \text{Gamma}
  \left(\frac{n}{2} - 1, \frac{y^T \left[A_{\gamma}^{-1} - A_{\gamma}^{-1} X (X^T
  A_{\gamma}^{-1} X)^{-1} X^T A_{\gamma}^{-1} \right] y}{2} \right).  
\end{align*}
Finally if we take a prior $p(\gamma)$ for $\gamma$, then the
posterior of $\gamma$ becomes:
\begin{align*}
  p(\gamma \mid y) \propto \frac{p(\gamma) |A_{\gamma}|^{-1/2} |X^T
  A_{\gamma}^{-1} X|^{-1/2}}{\left(y^T \left[A_{\gamma}^{-1} - A_{\gamma}^{-1} X (X^T
  A_{\gamma}^{-1} X)^{-1} X^T A_{\gamma}^{-1} \right] y \right)^{(n/2)
  - 1}}.   
\end{align*}

\section{Proof of the $C \rightarrow \infty$ fact}

\begin{fact}\label{Climit}
  The quantity
\begin{align*}
   \hat{f}(x) &= \left(C (1, x) X^T + \tau^2 \mathbf{k}_I^T \right)
                \left(\Sigma^{-1} - \Sigma^{-1} X \left(C^{-1} 
                            I_2 + X^T \Sigma^{-1} X  \right)^{-1} X^T
                \Sigma^{-1} \right) y
\end{align*}
converges, as $C \rightarrow \infty$, to:
\begin{align*}
   \hat{f}(x) := (1, x) \left(X^T A_{\gamma}^{-1} X \right)^{-1} X^T
  A_{\gamma}^{-1} y + \gamma^2 \mathbf{k}_I^T \left(A_{\gamma}^{-1} -
  A_{\gamma}^{-1} X (X^T A_{\gamma}^{-1} X)^{-1} X^T A_{\gamma}^{-1}
  \right) y. 
\end{align*}
Here $\tau = \gamma \sigma$ and $\Sigma = \tau^2 K_I + \sigma^2 I_n$.  
\begin{align*}
  \Sigma &= \tau^2 K_I + \sigma^2 I_n = \sigma^2 \gamma^2 K_I +
           \sigma^2 I_n = \sigma^2 A_{\gamma} 
\end{align*}
where $A_{\gamma} = I + \gamma^2 K_I$.  
\end{fact}

\begin{proof}[Proof of Fact \ref{Climit}]
  First note that
\begin{align*}
  \hat{f}(x)  &= \left(C (1, x) X^T + \tau^2 \mathbf{k}_I^T \right)
                \left(\Sigma^{-1} - \Sigma^{-1} X \left(C^{-1} 
                            I_2 + X^T \Sigma^{-1} X  \right)^{-1} X^T
                \Sigma^{-1} \right) y \\
  &= T_1 + T_2
\end{align*}
where
\begin{align*}
  T_1 := C (1, x) X^T\left(\Sigma^{-1} - \Sigma^{-1} X \left(C^{-1} 
                            I_2 + X^T \Sigma^{-1} X  \right)^{-1} X^T
                \Sigma^{-1} \right) y
\end{align*}
and
\begin{align*}
  T_2 &:= \tau^2 \mathbf{k}_I^T 
                \left(\Sigma^{-1} - \Sigma^{-1} X \left(C^{-1} 
                            I_2 + X^T \Sigma^{-1} X  \right)^{-1} X^T
        \Sigma^{-1} \right) y \\
  &\rightarrow \tau^2 \mathbf{k}_I^T 
                \left(\Sigma^{-1} - \Sigma^{-1} X \left(X^T
    \Sigma^{-1} X  \right)^{-1} X^T 
    \Sigma^{-1} \right) y ~~ \text{ as $C \rightarrow \infty$} \\
  &= \tau^2 \mathbf{k}_I^T \left(\sigma^{-2} A_{\gamma}^{-1} -
    \sigma^{-2}A_{\gamma}^{-1} X \left(X^T 
    A_{\gamma}^{-1} X  \right)^{-1} X^T 
    A_{\gamma}^{-1} \right) y \\
  &= \gamma^2 \mathbf{k}_I^T \left(A_{\gamma}^{-1} -
A_{\gamma}^{-1} X \left(X^T 
    A_{\gamma}^{-1} X  \right)^{-1} X^T 
    A_{\gamma}^{-1} \right) y. 
\end{align*}
The limit of $T_1$ as $C \rightarrow \infty$ is calculated as
follows. First note that
\begin{align*}
  T_1 &= C (1, x) X^T\left(\Sigma^{-1} - \Sigma^{-1} X \left(C^{-1} 
                            I_2 + X^T \Sigma^{-1} X  \right)^{-1} X^T
        \Sigma^{-1} \right) y \\
  &= C(1, x) X^T \Sigma^{-1} y - C (1, x) X^T \Sigma^{-1} X\left(C^{-1} 
                            I_2 + X^T \Sigma^{-1} X  \right)^{-1} X^T
    \Sigma^{-1} y \\
  &= C(1, x) X^T \Sigma^{-1} y - C(1, x) \left[\left(C^{-1} I_2 + X^T
    \Sigma^{-1} X \right) (X^T \Sigma^{-1} X)^{-1} \right]^{-1}   X^T
    \Sigma^{-1} y \\
  &= C(1, x) X^T \Sigma^{-1} y - C(1, x) \left[I_2 + C^{-1} (X^T
    \Sigma^{-1} X)^{-1}  \right]^{-1} X^T \Sigma^{-1} y
\end{align*}
Using the fact:
\begin{align*}
  \left(I + A^{-1}/C\right)^{-1} = I - A^{-1}/C + O(1/C^2)
\end{align*}
for $A = X^T \Sigma^{-1} X$, we obtain
\begin{align*}
  T_1 &= C(1, x) X^T \Sigma^{-1} y - C(1, x) \left(I - (X^T \Sigma^{-1}
        X)^{-1}/C +
        O(1/C^2) \right) X^T \Sigma^{-1} y \\
  &= (1, x) (X^T \Sigma^{-1} X)^{-1} (X^T \Sigma^{-1} y) + (1, x)
    O(1/C) X^T \Sigma^{-1} y \\
  &\rightarrow (1, x) (X^T \Sigma^{-1} X)^{-1} (X^T \Sigma^{-1} y)
    \text{ as $C \rightarrow \infty$} \\
  &= (1, x) (X^TA_{\gamma}^{-1} X)^{-1} (X^T A_{\gamma}^{-1} y)
    \text{ because $\Sigma = \sigma^2 A_{\gamma}$}. 
\end{align*}
This completes the proof of Fact \ref{Climit}. 
\end{proof} 






%\begin{thebibliography}{99}

%\bibitem[MacKay and Peto(1995)]{MacKay1995HierarchicalDirichlet}
%D.~J.~C. MacKay and L.~C. Peto,
%\newblock ``A hierarchical Dirichlet language model,''
%\newblock \emph{Natural Language Engineering}, vol.~1,
%no.~3, pp.~289--308, 1995.

%\bibitem[Rubin(1981)]{Rubin1981BayesianBootstrap}
%D.~B. Rubin,
%\newblock ``The Bayesian bootstrap,''
%\newblock \emph{The Annals of Statistics}, vol.~9,
%no.~1, pp.~130--134, 1981.
  
%\bibitem[Fisher(1934)]{Fisher1934}
%R.~A. Fisher,
%\newblock ``Probability, likelihood and quantity of information in the logic of
%uncertain inference,''
%\newblock \emph{Proceedings of the Royal Society of London. Series~A,
%Containing Papers of a Mathematical and Physical Character}, vol.~146,
%no.~856, pp.~1--8, 1934.

%\bibitem[Efron(2010)]{Efron}
%Efron, B. (2010)
%\textit{Large-scale inference: empirical Bayes methods for estimation,
%testing, and prediction}. Cambridge University Press, 2010.

%\bibitem[Shumway and Stoffer(2017)]{ShumwayStoffer}
%Shumway, R. H. and Stoffer, D. S. (2017)
%\textit{Time Series Analysis and Its Applications: With R
%  Examples}. Springer, 4th edition. 
  
%     \bibitem[Jaynes(2003)]{Jaynes} Jaynes, E. T, (2003)
%  \textit{Probability theory: the logic of science}. Cambridge
%  University Press, 2003.
 
%  \bibitem[Sinai(1992)]{Sinai} Sinai, Y. G, (1992)
%  \textit{Probability theory: an introductory course}. Springer
%  Verlag, 1992.  

%\bibitem[{deGroot(1986)}]{DeGrootBlackwell}DeGroot, M. H. (1986)
%       A conversation with David Blackwell.
%\newblock \emph{Statistical Science}, \textbf{1}, 40--53.

%         \bibitem[Mosteller(1987)]{Mosteller} Mosteller, F, (1987)
%  \textit{Fifty challenging problems in probability with
%    solutions}. Courier Corporation, 1987.

%    \bibitem[MacKay(2003)]{MacKay} MacKay, D. J. C., (2003)
%  \textit{Information theory, inference, and learning
%  algorithms}. Cambridge University Press, 2003.

%\bibitem[{Lindley(2013)}]{Lindley}Lindley, D. V. (2013)
%   \textit{Understanding Uncertainty}. John Wiley and sons, 2013.


%\end{thebibliography}







\end{document}

%%% Local Variables:
%%% mode: latex
%%% TeX-master: t
%%% End:
